% Auto-generated by build/rh_branch_successor.py. Do not edit by hand.

\section*{RH Branch Successor Pilot}

\addcontentsline{toc}{section}{RH Branch Successor Pilot}



\textbf{This is the highest-sensitivity governance pilot.} Any version of this packet that closes, promotes, or weakens \texttt{ob:RH.S1-ARC} from status~$[\mathrm{O}]$ would implicitly claim to prove the Riemann Hypothesis. The quality gate and equivalence check both enforce this.



\subsection*{Posture and S1 decomposition}

Computed RH posture: \textbf{\texttt{cert-proj}}. CERT-PROJ means: the branch architecture is constituted and the primary open obligation is Witness Assembly~S1, reduced to \texttt{ob:RH.S1-ARC} alone. The branch does \emph{not} claim to be near proving RH.



\begin{center}

\begin{tabular}{lll}

\textbf{S1 sub-predicate} & \textbf{Status} & \textbf{Route}\\

\hline

S1-TRACE & [C] & Scaling-flow candidate (conditional)\\

S1-DESCENT & [C] & Transport-localization chain L1--L7 (conditional)\\

S1-COMP & [C] & SCC-MCS discharged (SCC Branch)\\

\textbf{S1-ARC} & \textbf{[O]} & \textbf{IRREDUCIBLE --- mathematical content of RH}\\

\hline

\end{tabular}

\end{center}



\subsection*{Non-claims (strictly maintained)}

\begin{itemize}

\item The RH pilot does \textbf{not} prove the Riemann Hypothesis

\item S1-ARC is \textbf{not} conditionally discharged

\item S1-TRACE, S1-DESCENT, S1-COMP are [C], \textbf{not} [U]

\item No sub-predicate is promoted beyond its declared status

\item No arithmetic primitive is imported outside the declared branch structure

\end{itemize}



\begin{longtable}{p{0.30\linewidth}p{0.14\linewidth}p{0.10\linewidth}p{0.12\linewidth}p{0.20\linewidth}}
\caption{RH successor-pilot claim catalog}\label{tab:rh-successor-claims}\\
\textbf{ID} & \textbf{Kind} & \textbf{D.} & \textbf{C.} & \textbf{Scope}\\
\hline\endfirsthead
\textbf{ID} & \textbf{Kind} & \textbf{D.} & \textbf{C.} & \textbf{Scope}\\
\hline\endhead
thm:RH.1 & theorem & C & C & branch\\
thm:RH.2 & theorem & C & C & branch\\
thm:RH.3 & theorem & C & C & branch\\
ob:RH.4 & obligation & C & C & conditional-scaling-flow\\
ob:RH.5 & obligation & C & C & conditional-SCC-admissibility\\
ob:RH.6 & obligation & C & C & conditional-structural\\
ob:RH.S1-TRACE & obligation & C & C & conditional-scaling-flow\\
ob:RH.S1-DESCENT & obligation & C & C & conditional-L1-L7\\
ob:RH.S1-COMP & obligation & C & C & conditional-SCC-MCS\\
ob:RH.S1-ARC & frontier & O & O & arithmetic-irreducible\\
thm:RH.S1.formal & frontier & O & O & arithmetic-irreducible\\
status:RH & status\_note & cert-proj & cert-proj & branch\\
\end{longtable}



\subsection*{Generated claim shells}

\begin{nfctheorem}{RH1 — Canonical Object Xi}
\textbf{ID.} \texttt{thm:RH.1}\quad\textbf{Decl.} C\quad\textbf{Comp.} C\\
\textbf{Scope.} regime=RH-canonical-object; certificate\_scope=branch; promotion\_ceiling=C

The RH canonical object is Xi(t) = xi(1/2 + it), where xi is the completed Riemann xi-function. RH is the condition that all zeros of Xi are real, equivalently that all non-trivial zeros of the Riemann zeta function lie on the critical line Re(s) = 1/2. The branch target is established by this canonical identification: the NFC branch operates on Xi as its declared source object, descended from U via the SCC collar-algebra.

\medskip
\noindent\textit{Dependencies.} thm:SCC.MCS

\medskip
\noindent\textit{Proof stub.} Immediate from the definition of the RH canonical object in the branch-specific arena. The xi-function is the canonical completion of the Riemann zeta function, satisfying the functional equation xi(s) = xi(1-s). The branch target (RH: all zeros of Xi are real) is the canonical formulation adopted by the branch. No hypotheses beyond the branch architecture declaration are needed.
\end{nfctheorem}



\begin{nfctheorem}{RH2 — Spine Reduction}
\textbf{ID.} \texttt{thm:RH.2}\quad\textbf{Decl.} C\quad\textbf{Comp.} C\\
\textbf{Scope.} regime=RH-canonical-object; certificate\_scope=branch; promotion\_ceiling=C

Prime data (prime-power explicit formula), explicit-formula data (log-derivative trace formula), and zero data (the zero set of Xi) all descend from the canonical object O\_RH = Xi. The spine reduction establishes that the branch has a single canonical source object and that all arithmetic data relevant to RH is derivable from Xi without importing external number-theoretic primitives as independent assumptions.

\medskip
\noindent\textit{Dependencies.} thm:RH.1

\medskip
\noindent\textit{Proof stub.} By the branch architecture: the prime-power explicit formula is the log-derivative of Xi; the zero data is the zero set of Xi; the explicit-formula data descends via the Hadamard product representation. All three descend from O\_RH = Xi by declared derivation steps, without new source objects. The SCC-MCS carrier structure (thm:SCC.MCS) provides the minimal faithful carrier for this descent.
\end{nfctheorem}



\begin{nfctheorem}{RH3 — Admissible Probe Space}
\textbf{ID.} \texttt{thm:RH.3}\quad\textbf{Decl.} C\quad\textbf{Comp.} C\\
\textbf{Scope.} regime=RH-SCC-admissible; certificate\_scope=branch; promotion\_ceiling=C; hyp=SCC-M1-M4,SCC-RH-kernel-positivity-equivalence

The SCC-admissible probe space H\_RH is the Hilbert space of functions psi satisfying the four SCC admissibility conditions: Mellin structure SCC-M1, multiplicative convolution stability SCC-M2, Q\textasciicircum{}x-quotient descent SCC-M3, and L\textasciicircum{}2 membership SCC-M4. The RH kernel K(t,u) is positive-definite on H\_RH if and only if RH holds (SCC-RH kernel positivity equivalence). H\_RH is the natural Hilbert-Polya space for the RH branch.

\medskip
\noindent\textit{Hypotheses.} SCC-M1-M4, SCC-RH-kernel-positivity-equivalence

\medskip
\noindent\textit{Dependencies.} thm:RH.2, thm:SCC.MCS

\medskip
\noindent\textit{Proof stub.} By the SCC-RH kernel positivity equivalence (SCC Branch): RH iff K(t,u) positive-definite iff Q[psi] >= 0 for all SCC-admissible psi. The SCC-admissible probe space is therefore the natural Hilbert space for the RH kernel test. The space H\_RH with the L\textasciicircum{}2 inner product from SCC-M4 is the candidate Hilbert-Polya space consistent with necessary conditions RH-N1 through RH-N5.
\end{nfctheorem}



\begin{nfcobligation}{RH4 — Operator Realization (Scaling-Flow Candidate)}
\textbf{ID.} \texttt{ob:RH.4}\quad\textbf{Decl.} C\quad\textbf{Comp.} C\\
\textbf{Scope.} regime=RH-scaling-flow; certificate\_scope=conditional-scaling-flow; promotion\_ceiling=C; hyp=Phase-A,scaling-flow-candidate-Gsf

RH4 (Operator Realization): construct an explicit operator H\_RH satisfying at least one of three options: (i) Xi is the regularized spectral determinant det\_reg(H\_RH\textasciicircum{}2 - t\textasciicircum{}2); (ii) the trace formula reproduces the von Mangoldt explicit formula Tr\_\{G\_sf\}(h) = sum\_\{p,m\} (log p) p\textasciicircum{}\{-m/2\} h(m log p); (iii) the admissible eigenmodes correspond exactly to zeros of zeta.

Status: conditionally advanced by the scaling-flow candidate G\_sf and the orbit grammar program. Option (ii) is the primary realized route; it is conditionally realized by the scaling-flow candidate and the log-derivative trace formula (conditional [C]). Remaining: full branch-internal operator realization at canonical theorem force — logdiff-channel Clause 2 and trace-pairing law at unconditional level.

\medskip
\noindent\textit{Hypotheses.} Phase-A, scaling-flow-candidate-Gsf

\medskip
\noindent\textit{Dependencies.} thm:RH.3
\end{nfcobligation}



\begin{nfcobligation}{RH5 — Selection Exclusion (SCC Support-Rigidity)}
\textbf{ID.} \texttt{ob:RH.5}\quad\textbf{Decl.} C\quad\textbf{Comp.} C\\
\textbf{Scope.} regime=RH-SCC-admissible; certificate\_scope=conditional-SCC-admissibility; promotion\_ceiling=C; hyp=SCC-support-rigidity

RH5 (Selection Exclusion): prove that non-critical-line zeros correspond to non-admissible modes. The SCC support-rigidity chain (SCC Branch, thm:rh-ac-a3) shows that any off-critical-line zero fails the branch admissibility conditions in W\_RH. Status: conditionally hardened — the SCC support-rigidity provides the structural exclusion mechanism at SCC admissibility level. Remaining: lifting from SCC-admissibility to full operator-level exclusion once H\_RH is explicitly constructed (ob:RH.4).

\medskip
\noindent\textit{Hypotheses.} SCC-support-rigidity

\medskip
\noindent\textit{Dependencies.} thm:RH.3, thm:SCC.MCS
\end{nfcobligation}



\begin{nfcobligation}{RH6 — Critical-Line Localization}
\textbf{ID.} \texttt{ob:RH.6}\quad\textbf{Decl.} C\quad\textbf{Comp.} C\\
\textbf{Scope.} regime=RH-SCC-admissible; certificate\_scope=conditional-structural; promotion\_ceiling=C; hyp=Phase-A,SCC-RH-equivalence

RH6 (Critical-Line Localization): conditional on RH4 and RH5, prove that no SCC-admissible probe gives Q[psi] < 0 and therefore all admissible zeros satisfy Re(rho) = 1/2. Status: conditionally hardened by prop:rh-critical-line-structural [C] — the structural critical-line argument shows that RH4 + RH5 conditionally imply Q[psi] >= 0 for all admissible psi, giving RH at conditional force. Remaining gap to unconditional RH: full branch-internal proof of RH4 at canonical theorem force plus extension of the SCC/RH equivalence beyond Phase-A scope.

\medskip
\noindent\textit{Hypotheses.} Phase-A, SCC-RH-equivalence

\medskip
\noindent\textit{Dependencies.} ob:RH.4, ob:RH.5
\end{nfcobligation}



\begin{nfcobligation}{S1-TRACE — Trace Formula Pairing (Conditionally Discharged)}
\textbf{ID.} \texttt{ob:RH.S1-TRACE}\quad\textbf{Decl.} C\quad\textbf{Comp.} C\\
\textbf{Scope.} regime=RH-scaling-flow; certificate\_scope=conditional-scaling-flow; promotion\_ceiling=C; hyp=Phase-A,scaling-flow-candidate-Gsf

S1-TRACE: the trace formula pairing Tr\_\{G\_sf\}(h) = sum\_\{p,m\} (log p) p\textasciicircum{}\{-m/2\} h(m log p) is certified for all h in W\_RH. Status: CONDITIONALLY DISCHARGED. The scaling-flow candidate G\_sf conditionally realizes the log-derivative trace formula (RH-N2), giving S1-TRACE at conditional [C] force under the Phase-A and scaling-flow hypotheses. This sub-predicate does not require the arithmetic content of S1-ARC.

\medskip
\noindent\textit{Hypotheses.} Phase-A, scaling-flow-candidate-Gsf

\medskip
\noindent\textit{Dependencies.} ob:RH.4
\end{nfcobligation}



\begin{nfcobligation}{S1-DESCENT — Arithmetic Content Descends Without Hidden Import (Conditionally Discharged)}
\textbf{ID.} \texttt{ob:RH.S1-DESCENT}\quad\textbf{Decl.} C\quad\textbf{Comp.} C\\
\textbf{Scope.} regime=RH-transport-localization; certificate\_scope=conditional-L1-L7; promotion\_ceiling=C; hyp=transport-localization-chain-L1-L7,RH-defect-ledger

S1-DESCENT: the arithmetic content descends to the Mellin/Dirichlet arithmetic data for zeta(s) without hidden import. Status: CONDITIONALLY DISCHARGED via the transport-localization chain L1-L7 and the RH branch defect ledger D\_RH. The descent from Xi to the Mellin/Dirichlet arithmetic data is established conditionally. No hidden arithmetic primitive is imported outside the declared branch structure.

\medskip
\noindent\textit{Hypotheses.} transport-localization-chain-L1-L7, RH-defect-ledger

\medskip
\noindent\textit{Dependencies.} thm:RH.3
\end{nfcobligation}



\begin{nfcobligation}{S1-COMP — SCC-Minimal Faithful Witness Package (Discharged via SCC-MCS)}
\textbf{ID.} \texttt{ob:RH.S1-COMP}\quad\textbf{Decl.} C\quad\textbf{Comp.} C\\
\textbf{Scope.} regime=RH-SCC-admissible; certificate\_scope=conditional-SCC-MCS; promotion\_ceiling=C; hyp=SCC-MCS-discharged

S1-COMP: the arithmetic witness package is SCC-minimal faithful — no non-load-bearing arithmetic sector. Status: CONDITIONALLY DISCHARGED via SCC-MCS (SCC Branch, thm:SCC.MCS). The minimal faithful SCC carrier for the RH witness package exists and is uniquely determined by the SCC-MCS structure. The SCC-MCS discharge provides S1-COMP at the same conditional force as the SCC branch.

\medskip
\noindent\textit{Hypotheses.} SCC-MCS-discharged

\medskip
\noindent\textit{Dependencies.} thm:SCC.MCS
\end{nfcobligation}



\begin{nfcfrontier}{S1-ARC — Prove Q[psi]>=0 from Arithmetic Structure of Xi (IRREDUCIBLE FRONTIER)}
\noindent\colorbox{yellow!30}{\textbf{IRREDUCIBLE FRONTIER}} \textbf{ID.} \texttt{ob:RH.S1-ARC}\quad\textbf{Decl.} O\quad\textbf{Comp.} O\\
\textbf{Scope.} regime=RH-pure-arithmetic; certificate\_scope=arithmetic-irreducible; promotion\_ceiling=O

S1-ARC (IRREDUCIBLE ARITHMETIC FRONTIER): prove Q[psi] >= 0 for all psi in W\_RH directly from the arithmetic structure of Xi, WITHOUT assuming the Riemann Hypothesis.

This is the mathematical content of RH itself. It is NOT conditionally discharged. It is NOT approached by S1-TRACE, S1-DESCENT, or S1-COMP — those sub-predicates concern the structural machinery, not the positivity. S1-ARC is PURE ARITHMETIC: the question is whether the Mellin/Dirichlet descent structure of Xi forces the quadratic form Q to be non-negative, and this is precisely what the Riemann Hypothesis says.

Admissible discharge modes: explicit proof of Q[psi] >= 0 from the arithmetic structure of Xi; explicit supersession with proof of equivalence. Inadmissible: renaming the problem, showing Q[psi] >= 0 for a proper subset of W\_RH and calling it complete, or conditional arguments that assume RH.

Status: OPEN [O]. This is the irreducible remaining RH burden in the NFC canonical framework. The canon's current toolkit (collar-algebra, NFC spine, SCC machinery) does not provide the tools to resolve S1-ARC.
\end{nfcfrontier}



\begin{nfcfrontier}{RH S1 Formal Witness Assembly — Containing Frontier (Reduces to S1-ARC)}
\textbf{ID.} \texttt{thm:RH.S1.formal}\quad\textbf{Decl.} O\quad\textbf{Comp.} O\\
\textbf{Scope.} regime=RH-pure-arithmetic; certificate\_scope=arithmetic-irreducible; promotion\_ceiling=O

S1 (Witness Assembly) — Containing Frontier: the full Witness Assembly obligation reduces to S1-ARC alone. S1-TRACE, S1-DESCENT, and S1-COMP are all conditionally discharged; the irreducible remaining content is ob:RH.S1-ARC: proving Q[psi] >= 0 from the arithmetic structure of Xi without assuming RH. The S1 formal packet transforms "S1 is the irreducible arithmetic frontier" into a precise theorem statement with named sub-predicates, four of which are identified and three conditionally discharged. The residual is precisely the mathematical content of the Riemann Hypothesis.

\medskip
\noindent\textit{Dependencies.} ob:RH.S1-ARC
\end{nfcfrontier}



\begin{nfcstatusnote}{RH Branch Posture}
\textbf{ID.} \texttt{status:RH}\quad\textbf{Decl.} cert-proj\quad\textbf{Comp.} cert-proj\\
\textbf{Scope.} regime=RH-arithmetic-frontier; certificate\_scope=branch

RH branch posture: CERT-PROJ. The branch architecture is constituted; the primary open obligation is Witness Assembly S1, reduced to S1-ARC.

CERT-PROJ means: the branch has established the canonical object, probe space, and structural machinery for RH, and has conditionally discharged all obstacles except the irreducible arithmetic frontier S1-ARC. The branch is "certified to project" in the sense that it has a well-defined, sharply-stated remaining burden — but it does not claim to be near closure.

Active frontier: ob:RH.S1-ARC (Q[psi] >= 0 from arithmetic structure of Xi). All other S1 sub-predicates are [C]. All RH ladder obligations are [C] or conditionally advanced. 9 of 10 AC obstructions discharged.

Non-claims: the branch does not prove RH, does not claim to be near proving RH, and does not promote S1-ARC beyond [O] status.

\medskip
\noindent\textit{Dependencies.} thm:RH.3, ob:RH.4, ob:RH.5, thm:V.UBLT, thm:VII.6.4
\end{nfcstatusnote}


